% ===================================================================
%  Задача 4.4, п. 1: сдвиг функции потребления при переходе
%  от города к селу. Наклон (предельная склонность) общий,
%  свободный член отличается на gamma.
%      c_t = alpha + beta*y_t + gamma*r_t + eps_t
%  Компилировать: xelatex problem_4.4_shift.tex
% ===================================================================
\documentclass[border=6pt]{standalone}
% ===================================================================
%  Общая преамбула для отдельных картинок семинара 3.
%  Шрифты и цвета те же, что в ../../_tex/econ_preamble.tex.
%  Компилировать XeLaTeX'ом.
% ===================================================================
\usepackage{fontspec}
\setmainfont{Times New Roman}
\usepackage[russian]{babel}
\usepackage{amsmath,amssymb}
\usepackage{xcolor}
\usepackage{tikz}

\definecolor{accent}{RGB}{0,80,130}
\definecolor{solbar}{RGB}{0,120,90}
\definecolor{notebar}{RGB}{170,120,0}


\begin{document}
\begin{tikzpicture}[>=stealth, line join=round, x=0.9cm, y=0.9cm]
  % параметры картинки: alpha = 1, beta = 0.5, gamma = 1.2, y0 = 3
  \draw[->] (0,0) -- (8.6,0) node[anchor=north] {$y$};
  \draw[->] (0,0) -- (0,6.6) node[anchor=east]  {$c$};
  \node[anchor=north east] at (0,0) {$0$};

  % две функции потребления: город (r = 0) и село (r = 1)
  \draw[very thick, accent]  (0,1.0) -- (8,5.0)
        node[anchor=west] {город, $r_t = 0$: \ $c = \alpha + \beta y$};
  \draw[very thick, solbar]  (0,2.2) -- (8,6.2)
        node[anchor=west] {село, $r_t = 1$: \ $c = (\alpha + \gamma) + \beta y$};

  % свободные члены и сдвиг gamma
  \fill[accent] (0,1.0) circle (1.5pt);
  \fill[solbar] (0,2.2) circle (1.5pt);
  \node[anchor=east, accent] at (0,1.0) {$\alpha$};
  \node[anchor=east, solbar] at (0,2.2) {$\alpha + \gamma$};
  \draw[<->, thick] (1.0,1.5) -- (1.0,2.7)
        node[midway, anchor=west] {$\gamma$};

  % одинаковый наклон: треугольники приращений (у села — над прямой,
  % у города — под ней, чтобы подписи не налезали на соседнюю линию)
  \draw[black!60] (4.5,4.45) -- (4.5,5.45) -- (6.5,5.45);
  \node[anchor=east,  black!70] at (4.5,4.95) {\small $\beta\,\Delta y$};
  \node[anchor=south, black!70] at (5.5,5.45) {\small $\Delta y$};
  \draw[black!60] (5.0,3.5) -- (7.0,3.5) -- (7.0,4.5);
  \node[anchor=north, black!70] at (6.0,3.5) {\small $\Delta y$};
  \node[anchor=west,  black!70] at (7.0,4.0) {\small $\beta\,\Delta y$};

  % средняя склонность при одном и том же доходе y0 — наклоны лучей из нуля
  \draw[dotted, black!55] (3,0) -- (3,3.7);
  \node[anchor=north] at (3,0) {$y_0$};
  \draw[dashed, accent]  (0,0) -- (3,2.5);
  \draw[dashed, solbar]  (0,0) -- (3,3.7);
  \fill[accent] (3,2.5) circle (1.5pt);
  \fill[solbar] (3,3.7) circle (1.5pt);

  % подпись
  \node[anchor=north west, align=left, font=\small\itshape, text width=13cm]
    at (-0.6,-0.9)
    {Предельная склонность $\partial c/\partial y = \beta$ одинакова
     (треугольники приращений равны). Средняя склонность $c/y$ — наклон
     пунктирного луча из нуля — при одном и том же доходе $y_0$ различается:
     $\beta + (\alpha+\gamma)/y_0$ у села против $\beta + \alpha/y_0$ у города.
     Проверка: $H_0\colon \gamma = 0$, $t$-тест с $n-3$ степенями свободы.};
\end{tikzpicture}
\end{document}
